% !PRINTED 285
\weberpage{285}
One thus finds, by an easy, if somewhat laborious, calculation
\weq{15}{\begin{array}{rl@{\qquad}rl}
\omega\tau &{}= \tau^2\chi^2\omega\Theta^2, & \Theta\tau &{}= \tau^3\omega\Theta\\
\omega\tau^2 &{}= \tau\chi^2\Theta^3, & \Theta\tau^2 &{}= \tau\chi\omega\Theta^3\\
\omega\tau^3 &{}= \tau^5\chi\Theta^2, & \Theta\tau^3 &{}= \tau^5\chi\omega\\
\omega\tau^4 &{}= \tau^4\chi^2\Theta, & \Theta\tau^4 &{}= \tau^6\omega\Theta^2\\
\omega\tau^5 &{}= \tau^3\chi^2\omega\Theta, & \Theta\tau^5 &{}= \tau^2\Theta^3\\
\omega\tau^6 &{}= \tau^6\omega\Theta^3, & \Theta\tau^6 &{}= \tau^4\chi^2\Theta^2.
\end{array}}
But all these twelve relations can be derived as consequences
from four of them, which may be selected in manifold ways.
One can, for example, choose for these four fundamental
relations the following:
\weq{16}{\begin{array}{c}
\omega\tau = \tau^2\chi^2\omega\Theta^2, \quad \omega\tau^3 = \tau^5\chi\Theta^2, \quad \omega\tau^4 = \tau^4\chi^2\Theta,\\
\Theta\tau = \tau^3\omega\Theta,
\end{array}}
from which, with the use of the formulae (3), (12), (13), all the others
follow easily, e.g.:
\[
\omega\tau^5 = \tau^4\chi^2\Theta\tau = \tau^4\chi^2\tau^3\omega\Theta = \tau^3\chi^2\omega\Theta,
\]
and so the rest.

Just as we have seen earlier that we can regard $\chi$ and $\Theta$
as generating elements of the octahedral group, so we can
now regard the two elements $\omega, \tau$ as generators of the group
$L_7$, for it follows from (15):
\weq{17}{\Theta^3 = \omega\tau\omega\tau^6, \quad \chi = \tau^2\Theta\tau^3\omega.}
Summing up, we obtain the following result:

\gesperrt{I. If four elements $\tau, \chi, \omega, \Theta$ of the degrees 7, 3, 2, 4
are given, in whose composition the relations
 hold:}
\[
\begin{array}{ccc}
\omega\chi = \chi^2\omega, & \Theta\omega = \omega\Theta^3, & \Theta\chi = \chi^2\omega\Theta^2,\\
\Theta^2\chi = \chi\omega\Theta^3, & \omega\tau = \tau^2\chi^2\omega\Theta^2, & \omega\tau^3 = \tau^5\chi\Theta^2,\\
\omega\tau^4 = \tau^4\chi^2\Theta, & \Theta\tau = \tau^3\omega\Theta, & \chi\tau = \tau^4\chi,
\end{array}
\]
\gesperrt{then the 168 elements}
\[
\sigma = \tau^\varrho \chi^\lambda \omega^u \Theta^v,
\]
\gesperrt{when $\varrho, \lambda, \mu, \nu$ run through complete systems of residues
 with respect to the moduli 7, 3, 2, 4, form a simple group
of degree 168, and $\omega$ and $\tau$ can be regarded as generating
 elements of this group.} % [sic] exponents in the enunciation text (mu, nu) do not match the exponents u, v used in the preceding formula for sigma
